% TOP_PROBLEM: {"id": 3200013, "title": "Exact Dedekind numbers", "category_id": 2, "category": {"id": 2, "name": "combinatorics", "display_name": "Combinatorics", "description": "Counting problems, graph theory, discrete structures.", "slug": "combinatorics", "order_index": 2, "created_at": "2026-07-31T15:26:25.671Z"}, "status": "open"}
% FIRST_POSTED: null
% ENTRY_KIND: statement_only
% Imported from: batch12.tex; UnsolvedMath ID 3200013
% Compile twice: lualatex 3200013.tex
\documentclass[11pt]{article}
\usepackage{fontspec}
\setmainfont{Latin Modern Roman}
\usepackage{amsmath,amssymb,mathtools,unicode-math}
\setmathfont{Latin Modern Math}
\usepackage{xurl,hyperref}
\hypersetup{hidelinks,pdftitle={UnsolvedMath Problem 3200013}}
\newcommand{\sref}[2]{\hyperlink{source:#1:#2}{[S#2]}}
\newcommand{\cataloguescope}[1]{\paragraph{Mathematical question.}#1}
\begin{document}
% UnsolvedMath ID 3200013; source catalogue code AMR-031-0013
\setcounter{section}{3200012}
\section{Exact Dedekind numbers}\label{3200013}
{\small\sffamily UnsolvedMath ID 3200013 · Combinatorics}
\subsection{Definitions and mathematical statement}

\paragraph{Shared notation.}$\mathbb N=\{1,2,\ldots\}$, and $\mathbb Z,\mathbb Q,\mathbb R,\mathbb C$ denote the integers, rationals, reals and complex numbers. Any other conventions are specified locally.


Write $[n]=\{1,\ldots,n\}$ and let $2^{[n]}$ be its power set. An antichain is a family $\mathcal A\subseteq2^{[n]}$ such that no distinct $A,B\in\mathcal A$ satisfy $A\subseteq B$. Define the Dedekind number
\[
M(n)=\#\{\mathcal A\subseteq2^{[n]}:\mathcal A\text{ is an antichain}\}.
\]
The empty family and the family $\{\varnothing\}$ are both counted. Equivalently, $M(n)$ counts all maps $f:\{0,1\}^n\to\{0,1\}$ satisfying $f(x)\le f(y)$ whenever $x_i\le y_i$ for every $i$. The correspondence associates to $f$ the inclusion-minimal subsets on which $f$ has value $1$; thus both constant functions are included. Counting antichains removes the harmless choice between increasing and decreasing conventions for Boolean functions.
\paragraph{Statement recorded in the supplied source.}
Determine the exact integer $M(n)$ for every $n\ge10$. The cited computation establishes $M(9)$ and develops exact counting formulae; an asymptotic estimate or a non-evaluated summation is not an exact value of the requested numbers. \sref{3200013}{2}
\subsection{Short English statement}
Count exactly the antichains in the subsets of an $n$-element set, or equivalently all monotone Boolean functions, for every $n\ge10$.
\subsection{Sources}
\begin{description}
\item[\hypertarget{source:3200013:1}{S1}] ULAM.AI and the UnsolvedMath Contributors, UnsolvedMath: A Curated Collection of Open Mathematics Problems, record 3200013 (AMR-031-0013). CC BY 4.0. \url{https://huggingface.co/datasets/ulamai/UnsolvedMath}.
\item[\hypertarget{source:3200013:2}{S2}] Patrick De Causmaecker and Lennart Van Hirtum, Solving systems of equations on antichains for the computation of the ninth Dedekind Number (2024), §§1,2.1–2.2, pp.1–5. \url{https://arxiv.org/abs/2405.20904}.
\end{description}
\label{end:3200013}
\end{document}
